\documentclass[conference]{IEEEtran}
\usepackage[T1]{fontenc}
\usepackage{graphicx,booktabs,amsmath,amssymb,xcolor,url,cite,multirow,array}
\usepackage[hidelinks]{hyperref}
\hyphenation{Align-er phone-me}

% ---------- helpers ----------
\newcommand{\tbd}[1]{\textcolor{red}{\textbf{[#1]}}}        % unfilled value
\newcommand{\figbox}[3]{%                                    % file, caption, label
  \begin{figure}[!t]\centering
  \IfFileExists{figures/#1}{\includegraphics[width=\columnwidth]{figures/#1}}%
  {\fbox{\parbox{0.9\columnwidth}{\centering\tbd{missing figure \detokenize{#1}; run notebook 09}}}}%
  \caption{#2}\label{fig:#3}\end{figure}}
\newcommand{\safeinput}[1]{\IfFileExists{#1}{\input{#1}}{\par\tbd{missing \detokenize{#1}; run notebook 09}\par}}
% ======================================================================
%  RESULT MACROS: the ONLY file you need to edit with numbers.
%  Every red [..] in the PDF is a macro you have not filled yet.
%  Replace \tbd{...} by the plain value, e.g. \newcommand{\nFailed}{37}
% ======================================================================
\newcommand{\repoURL}{https://github.com/Shreyashgol/phonetic-guided-speech-generation}

% ---- Filled from your executed notebooks 01-08 ----
\newcommand{\nClips}{13{,}100}
\newcommand{\totalHours}{23.92}
\newcommand{\minDur}{1.11}
\newcommand{\meanDur}{6.57}
\newcommand{\maxDur}{10.10}
\newcommand{\vocabSize}{14{,}058}
\newcommand{\nSingletons}{5{,}513}
\newcommand{\nRawNeqNorm}{1{,}505}
\newcommand{\nTrainIds}{11{,}790}   % 90\% of 13,100
\newcommand{\nValIds}{655}
\newcommand{\nTestIds}{655}

% ---- Notebook 03 (alignment + QA) ----
\newcommand{\mfaVersion}{3.4.3}
\newcommand{\nOOVtypes}{420}
\newcommand{\nOOVtokens}{1{,}863}
\newcommand{\pctOOVtokens}{0.83\%}
\newcommand{\nFailed}{0}
\newcommand{\nAligned}{13{,}100}
\newcommand{\nFlagSpn}{1{,}562}
\newcommand{\nFlagRate}{132}
\newcommand{\nFlagMin}{0}
\newcommand{\nFlagSil}{8}
\newcommand{\nExcluded}{1{,}684}
\newcommand{\pctExcluded}{12.9\%}
\newcommand{\nKept}{11{,}416}

% ---- Notebook 04 (EDA) ----
\newcommand{\nSpeechPhones}{777{,}055}
\newcommand{\meanPhoneFrames}{7.46}
\newcommand{\meanVowel}{7.69}
\newcommand{\meanStop}{7.33}
\newcommand{\stressGap}{3.56}
\newcommand{\pauseGap}{6.14}
\newcommand{\wordFinalGap}{0.59}

% ---- Notebook 06/07/08 (models) ----
\newcommand{\bestModel}{BiLSTM}
\newcommand{\maeBase}{2.400}
\newcommand{\maeBaseMs}{27.9}
\newcommand{\maeBest}{1.330}
\newcommand{\maeBestMs}{15.4}
\newcommand{\ciBestLow}{1.310}
\newcommand{\ciBestHigh}{1.353}
\newcommand{\gainFrames}{1.071}
\newcommand{\gainLow}{1.048}
\newcommand{\gainHigh}{1.092}
\newcommand{\rsqBest}{0.775}
\newcommand{\nnParams}{650{,}293}
\newcommand{\nnEpochs}{25}
\newcommand{\topAblation}{the identity features (phone, stress, class)}
\newcommand{\worstClass}{glides (relative MAE 20.4\%)}
\newcommand{\gpuName}{a Colab CUDA GPU}

% ---- Team (the viva is individual: fill truthfully) ----
\newcommand{\teamA}{Shreyash Golhani}
\newcommand{\teamB}{Vaageesh Kumar Singh}
\newcommand{\teamC}{Anwesha Adhikari}
\newcommand{\teamD}{Ananya Narang}
\newcommand{\teamE}{Aman Kumar}


\setlength{\textfloatsep}{10pt plus 2pt minus 4pt}
\setlength{\floatsep}{8pt plus 2pt minus 2pt}
\setlength{\intextsep}{8pt plus 2pt minus 2pt}
\renewcommand{\topfraction}{0.9}\renewcommand{\textfraction}{0.1}\renewcommand{\floatpagefraction}{0.8}
\begin{document}

\title{Phonetic Guided Speech Generation, Phase~1:\\Forced Alignment and Duration Prediction (ML Baseline)}
\author{\IEEEauthorblockN{\teamA, \teamB, \teamC, \teamD, \teamE}
Code: \url{\repoURL}}
\maketitle

% =====================================================================
\begin{abstract}
Text-to-speech (TTS) needs to know \emph{when} each phoneme occurs and \emph{how long} it lasts.
This report covers Phase~1 of a project comparing explicit phoneme-duration alignment with a learned acoustic generator.
Phase~1 is the classical machine-learning part: (i) forced alignment of LJSpeech with the Montreal Forced Aligner (MFA) to obtain phoneme-level ground truth,
(ii) exploratory analysis of the resulting durations, and (iii) a ladder of duration predictors from group-mean baselines to ridge regression, gradient-boosted trees and a BiLSTM.
\nFailed{} of \nClips{} utterances failed to align; \nExcluded{} (\pctExcluded{}) were removed by quality checks, mostly because of out-of-dictionary words.
On a held-out, utterance-disjoint test set the best model (\bestModel) reaches a mean absolute error of \maeBest{} frames (\maeBestMs{} ms),
against \maeBase{} frames for the per-phoneme-mean baseline; the improvement is \gainFrames{} frames (95\% bootstrap CI \gainLow--\gainHigh).
We also analyse where the models fail and which feature groups matter.
\end{abstract}

\begin{IEEEkeywords}
forced alignment, phoneme duration, text-to-speech, hidden Markov models, LJSpeech, Montreal Forced Aligner
\end{IEEEkeywords}

% =====================================================================
\section{Introduction}
\subsection{Problem and scope}
An explicit duration model tells a TTS system how many acoustic frames each phoneme occupies, but it does not produce timbre, coarticulation or prosody.
The full project (Phases~1 and 2) asks whether a learned acoustic generator fixes what duration-only synthesis lacks.
\textbf{Phase~1 covers only the ML component}: alignment, data analysis, feature engineering and duration regression. No audio is synthesised here.

\subsection{Objectives}
\begin{enumerate}
\item Produce reliable phoneme-level duration labels for LJSpeech and quantify their quality.
\item Characterise duration statistics and test them against the HMM duration model that underlies forced alignment.
\item Build and fairly compare duration predictors of increasing capacity.
\item Document failure modes and set up reusable assets for Phase~2.
\end{enumerate}

\subsection{Contributions}
\begin{itemize}
\item A reproducible pipeline: text cleaning $\rightarrow$ MFA alignment $\rightarrow$ QA $\rightarrow$ features $\rightarrow$ models $\rightarrow$ evaluation, driven by one shared audio configuration.
\item Alignment QA with four proxy checks that detect \emph{bad} (not only failed) alignments.
\item A frozen, utterance-level train/validation/test split to be reused unchanged in Phase~2.
\item Utterance-bootstrap confidence intervals, a paired significance test and a feature-group ablation.
\end{itemize}

\subsection{Evaluation map}
Table~\ref{tab:map} links each Phase~1 evaluation criterion to where it is addressed.

\begin{table}[t]\centering\caption{Where each Phase~1 criterion is addressed.}\label{tab:map}
\footnotesize
\begin{tabular}{@{}p{2.5cm}p{1.2cm}p{3.9cm}@{}}
\toprule
Criterion & Section & Evidence in repository \\ \midrule
Literature review & \ref{sec:lit} & references; MFA, FastSpeech, Grad-TTS \\
Dataset quality \& EDA & \ref{sec:data} & notebooks 01--04; \texttt{alignment\_qa.csv} \\
Feature engineering & \ref{sec:feat} & notebook 05; \texttt{feature\_sets.py} \\
Theoretical rigor & \ref{sec:theory} & derivations; HMM-vs-data test (nb.~04) \\
Model application & \ref{sec:models},\ref{sec:results} & notebooks 06--08 \\
Code quality & \ref{sec:repro} & public repo, configs, frozen splits \\
Presentation, viva & \ref{sec:contrib} & per-member contribution table \\ \bottomrule
\end{tabular}
\end{table}

% =====================================================================
\section{Literature Review}\label{sec:lit}
\textbf{Forced alignment.}
Classical alignment models each phoneme as a left-to-right hidden Markov model (HMM) with Gaussian-mixture emissions over spectral features and finds the best state path by Viterbi decoding \cite{rabiner1989,viterbi1967}; parameters are trained with expectation--maximisation (EM) \cite{dempster1977}.
The Montreal Forced Aligner (MFA) packages this on top of Kaldi \cite{povey2011}, training monophone then triphone models and offering pretrained English models and dictionaries \cite{mcauliffe2017}.
Alignment quality depends on the pronunciation dictionary and transcript fidelity, so out-of-vocabulary words and noisy transcripts are the main failure sources.

\textbf{Duration modelling in TTS.}
In HMM-based synthesis, duration is implicit in state self-transition probabilities, which yields geometric-like distributions that fit natural speech poorly \cite{rabiner1989}.
Neural TTS made duration explicit: FastSpeech trains a duration predictor on alignments from a teacher and uses a length regulator to expand phoneme encodings to frame length \cite{ren2019}; FastSpeech~2 predicts log-domain durations directly from forced-alignment targets \cite{ren2021}.
Our Phase~1 predictors follow the same recipe (log-domain targets, alignment-derived labels).

\textbf{Diffusion-based acoustic generation.}
Grad-TTS generates mel-spectrograms with a score-based diffusion decoder conditioned on text and a duration alignment \cite{popov2021}; a neural vocoder such as HiFi-GAN converts mels to waveforms \cite{kong2020}.
This is the Phase~2 context; Phase~1 supplies its duration labels and baseline.

\textbf{Gap addressed.}
Published systems rarely report how much duration is predictable from text alone, nor how alignment label noise limits the achievable error.
We measure both on a single-speaker corpus.

% =====================================================================
\section{Dataset, Alignment and EDA}\label{sec:data}
\subsection{Corpus}
LJSpeech 1.1 \cite{ljspeech2017}: \nClips{} clips from one female speaker reading non-fiction (public domain), $\approx$\totalHours{} h in total.
Clip length: min \minDur{} s, mean \meanDur{} s, max \maxDur{} s. All files are 22{,}050~Hz, mono, 16-bit PCM (checked on every file).
The transcript file has raw and normalised text; \nRawNeqNorm{} rows differ (numbers and abbreviations are expanded in the normalised column, which we use).
After our cleaning (lowercase, no punctuation) the text contains \vocabSize{} distinct words, of which \nSingletons{} occur once.

\subsection{Text preparation}
For MFA we lowercase, strip accents (Unicode NFKD), replace hyphens by spaces and remove all punctuation. Punctuation is not discarded: each word keeps a code (none / comma-like / sentence-final) used later as a feature.
A \emph{shared phoneme inventory} (15 vowels, 24 consonants in ARPAbet, plus \texttt{sil} and \texttt{spn}) is defined once and imported by both phases.

\subsection{Forced alignment}
We use MFA \mfaVersion{} with the pretrained \texttt{english\_us\_arpa} acoustic model and dictionary, single-speaker mode, no retraining; aligning all \nClips{} clips took about 22 minutes on a Colab CPU.
Phone boundaries are converted from seconds to mel frames at 22{,}050~Hz, hop 256 (11.6~ms/frame) with \emph{cumulative rounding}, so intervals remain contiguous. The shared audio configuration (FFT 1024, hop 256, 80 mels, 0--8~kHz) matches common pretrained vocoders.

\textbf{Dictionary coverage.} \nOOVtypes{} word types (\nOOVtokens{} tokens, \pctOOVtokens{} of all tokens) are absent from the dictionary; MFA maps them to the unknown-noise phone \texttt{spn}. The most frequent are abbreviations and names (\emph{mrs}, \emph{mr}, \emph{fbi}, \emph{motorcade}, \emph{connally}).

\textbf{Failures.} \nFailed{} of \nClips{} utterances produced no alignment and are logged in \texttt{failed\_utterances.csv}.

\textbf{Quality control for bad alignments.}
No ground truth exists, so four proxy flags are computed per utterance: (i) any \texttt{spn} phone (\nFlagSpn{} utterances); (ii) speech rate (phones per second of speech) outside the 0.5th--99.5th percentile, i.e.\ 8.3--14.3 phones/s (\nFlagRate{}); (iii) more than 20\% of speech phones at $\le$2 frames (\nFlagMin{}, so this check never triggered); (iv) leading silence $>1$~s or trailing silence $>1.5$~s (\nFlagSil{}).
The flags overlap; \nExcluded{} utterances (\pctExcluded{}) are excluded and \nKept{} remain, dominated by the \texttt{spn} flag. Random and flagged clips were also inspected visually (boundaries on the mel spectrogram) and by listening.

\subsection{Exploratory analysis}
The \nSpeechPhones{} retained speech phones have mean duration \meanPhoneFrames{} frames (median 7, maximum 93). 293 phones rounded to zero frames and are clipped to one frame.
\begin{itemize}
\item \textbf{Skew.} Durations are right-skewed; $\log d$ is far more symmetric (Fig.~\ref{fig:hist}), motivating log-domain regression.
\item \textbf{Class.} Class means range from 5.78 frames (glides) and 6.07 (nasals) to 8.27 (fricatives) and 10.39 (affricates); vowels average \meanVowel{} and stops \meanStop{}, but vowels vary most (std 4.66 frames). Individual phones differ more than classes: DH averages 3.67 frames and AW 15.31 (Fig.~\ref{fig:box}).
\item \textbf{Context.} Primary-stressed vowels are longer than unstressed ones by \stressGap{} frames (9.37 vs.\ 5.81); phones before a pause by \pauseGap{} frames (13.28 vs.\ 7.14); word-final phones by only \wordFinalGap{}. Interior pauses have a median length of 15 frames. These effects define the feature groups of Sec.~\ref{sec:feat}.
\end{itemize}
\figbox{overall_hist.png}{Duration of all speech phones in frames (left) and on log scale (right).}{hist}
\figbox{per_phone_box.png}{Duration per phoneme, sorted by mean.}{box}
\safeinput{tables/dataset.tex}

% =====================================================================
\section{Theoretical Foundations}\label{sec:theory}
\subsection{HMM forced alignment}
Each phoneme is a 3-state left-to-right HMM with transition probabilities $a_{ij}$ and emission density
$b_j(\mathbf{o})=\sum_{m}c_{jm}\,\mathcal{N}(\mathbf{o};\boldsymbol\mu_{jm},\boldsymbol\Sigma_{jm})$ over MFCC frames $\mathbf{o}_t$ (10~ms shift).
Given the transcript, the phoneme HMMs are concatenated into one utterance HMM, so only the state path is unknown. Viterbi decoding finds
\begin{equation}
\delta_t(j)=\max_i\big[\delta_{t-1}(i)\,a_{ij}\big]\,b_j(\mathbf{o}_t),
\end{equation}
and back-pointers recover the best path; phoneme boundaries are the points where the path leaves one phoneme's states.
Acoustic parameters are estimated by EM: the E-step uses forward--backward posteriors (soft state occupancies, summed over all paths), the M-step re-estimates means, covariances, mixture weights and transitions \cite{dempster1977,rabiner1989}.
\emph{Viterbi answers ``what is the single best path''; forward--backward answers ``how likely is state $s$ at frame $t$ over all paths''.}

\subsection{Duration under an HMM}
A state with self-loop probability $a$ gives a geometric duration $P(d)=a^{d-1}(1-a)$ with mean $1/(1-a)$, whose mode is always 1.
A phone of three left-to-right states has the sum of three such durations, a negative binomial shifted by 3:
\begin{equation}
P(d)=\binom{d-1}{2}(1-a)^3a^{\,d-3},\quad d\ge3,\quad \mathbb{E}[d]=\frac{3}{1-a}.
\end{equation}
Both have a fixed shape that cannot adapt to context; Sec.~\ref{sec:eda-theory} compares them to the data.

\subsection{Regression on log-duration}
Let $d_i\ge1$ be the duration in frames of phone $i$ and $y_i=\log d_i$. A predictor $f$ is trained by (masked) squared error in log space,
\begin{equation}
\mathcal{L}=\frac{1}{|\mathcal{M}|}\sum_{i\in\mathcal{M}}\big(f(\mathbf{x}_i)-y_i\big)^2,\qquad \hat d_i=\exp f(\mathbf{x}_i).
\end{equation}
Because $\log$ is monotone, $\exp f$ estimates the conditional median when log-residuals are roughly symmetric, and the median minimises expected absolute error, our main metric.

\subsection{Model classes}
\textbf{Ridge regression} minimises $\|\mathbf{y}-X\mathbf{w}\|^2+\alpha\|\mathbf{w}\|^2$; the $L_2$ penalty shrinks weights of rare one-hot categories \cite{hoerl1970}.
\textbf{Gradient-boosted trees} add trees fitted to the residual gradient of the loss \cite{friedman2001}, capturing non-linear interactions (e.g.\ stress $\times$ pause).
\textbf{BiLSTM} reads the whole phone sequence in both directions \cite{hochreiter1997}; right context matters because phrase-final lengthening depends on what follows.

% =====================================================================
\section{Feature Engineering}\label{sec:feat}
\textbf{Unit and target.} One row per speech phone; silence and \texttt{spn} rows are not modelled but signalled by \texttt{pause\_before/after}. Target: duration in frames (clipped at 1).
Table~\ref{tab:feat} lists the five feature groups.

\begin{table}[t]\centering\caption{Feature groups (identical for all models).}\label{tab:feat}
\footnotesize\begin{tabular}{@{}p{1.6cm}p{4.2cm}p{1.6cm}@{}}
\toprule
Group & Features & Motivation \\ \midrule
identity & phone, stress, class, is-vowel & intrinsic duration \\
context & previous/next two phones & coarticulation \\
word & position and length in word, first/last & word-final lengthening \\
utterance & word index, relative position, lengths & phrase effects \\
pause/punct. & pause before/after, punctuation after word & phrase-final lengthening \\ \bottomrule
\end{tabular}
\end{table}

\textbf{Splitting and leakage control.}
Phones of one sentence are strongly correlated, so splitting by phone would leak. We split by \emph{utterance} over all \nClips{} ids with a fixed seed (42): train/val/test = \nTrainIds{}/\nValIds{}/\nTestIds{} utterances (90/5/5\%), stored as files and never regenerated. Alignment exclusions are applied afterwards, so the test set is stable.
Hyper-parameters are tuned on validation only; the test set is scored once.

% =====================================================================
\section{Models and Experimental Protocol}\label{sec:models}
\subsection{Models}
\begin{itemize}
\item \textbf{Global mean}: one value (floor).
\item \textbf{Phone mean} (AML baseline): median duration per phoneme.
\item \textbf{Back-off table}: median per (phone, stress, pause-after), backing off to (phone, stress) then phone when fewer than 20 training rows.
\item \textbf{Ridge}: one-hot categoricals, standardised numerics, target $\log d$; $\alpha\in\{0.1,1,10,100\}$.
\item \textbf{Gradient-boosted trees} (histogram-based): categorical splits; (learning rate, iterations) $\in\{(0.1,200),(0.05,400)\}$.
\item \textbf{BiLSTM}: embeddings (phone 64, stress 8, class 8, punctuation 4) and a linear layer on 11 numeric/flag inputs, 2-layer bidirectional LSTM (hidden 128, dropout 0.2), MLP $256\!\to\!64\!\to\!1$; \nnParams{} parameters. AdamW (lr $2\!\times\!10^{-3}$, weight decay $10^{-4}$), batch 64, $\le$30 epochs, learning rate halved on plateau, gradient clipping 1.0, early stopping (patience 5) on validation MAE; \nnEpochs{} epochs run; seed 0; trained on \gpuName{}.
\end{itemize}

\subsection{Metrics}
With $\hat d_i$ the prediction and $N$ test phones: $\mathrm{MAE}=\frac1N\sum|\hat d_i-d_i|$, $\mathrm{RMSE}=\sqrt{\frac1N\sum(\hat d_i-d_i)^2}$, $R^2=1-\sum(\hat d_i-d_i)^2/\sum(d_i-\bar d)^2$. Frames convert to ms by $\times11.6$.
\textbf{Confidence intervals}: 1000 bootstrap resamples of \emph{utterances} (not phones, which are not independent); 95\% percentile intervals. \textbf{Significance}: paired bootstrap of the MAE difference between baseline and best model on identical resamples.

% =====================================================================
\section{Results}\label{sec:results}
\subsection{Model comparison}
\safeinput{tables/models.tex}
\figbox{models.png}{Test MAE with 95\% utterance-bootstrap intervals.}{models}
Table~\ref{tab:models} shows the test results. The best model, \bestModel, has MAE \maeBest{} frames (\maeBestMs{} ms; CI \ciBestLow--\ciBestHigh) and $R^2=\rsqBest$.
Relative to the phone-mean baseline (\maeBase{} frames, \maeBaseMs{} ms), the paired MAE gain is \gainFrames{} frames (95\% CI \gainLow--\gainHigh).
The ranking is monotone in model capacity: global mean 3.138 $>$ phone mean 2.400 $>$ back-off 2.220 $>$ Ridge 2.013 $>$ gradient boosting 1.540 $>$ BiLSTM 1.330 frames. The BiLSTM beats gradient boosting by 0.21 frames (1.330 vs.\ 1.540; $R^2$ 0.775 vs.\ 0.697) and the 95\% intervals ([1.310, 1.353] vs.\ [1.517, 1.566]) do not overlap. Ridge's validation MAE was insensitive to $\alpha$ (2.021--2.022), so its limit is linearity, not regularisation.

\subsection{Ablation}
\safeinput{tables/ablation.tex}
Removing \topAblation{} hurts most (Table~\ref{tab:ablation}). Context features come next (+0.402 frames), then pause/punctuation (+0.114) and word features (+0.061); utterance-level features give no measurable benefit ($-0.006$). Identity features alone reach 2.388 frames, about the phone-mean baseline (2.400), so context, pause and word information account for 0.85 frames of improvement. This matches Sec.~\ref{sec:data}: the pre-pause effect (+6.14 frames) is large, and only 5\% of phones precede a pause, yet \texttt{pause\_after} still matters because those phones carry the largest errors (error analysis below).
\figbox{ablation.png}{Increase in test MAE when a feature group is removed.}{abl}

\subsection{Error analysis}
\figbox{error_analysis.png}{Best model: predicted vs.\ true duration (left) and MAE by true-duration bin (right).}{err}
The largest relative error is in \worstClass{} (Fig.~\ref{fig:per_phone}). Errors show \emph{shrinkage toward the mean}: phones of $\le3$ frames are over-predicted by 0.88 frames on average, while phones of 12--20 and $>20$ frames are under-predicted by 1.77 and 4.79 frames (MAE 2.52 and 5.13). Phones followed by a pause have MAE 2.64 against 1.26 otherwise. By absolute error the diphthongs and long vowels are worst (OY 2.21, AY 2.15, AE 2.11, AW 2.05 frames); by relative error glides (0.204), nasals (0.199) and liquids (0.193) are worst, with vowels at 0.184.
The worst utterances (LJ024-0071, LJ017-0250, LJ036-0073) have mean errors of 3.17, 2.33 and 2.10 frames against 1.33 overall; Six of the ten worst utterances come from only two chapters: LJ028 (four; archaic wording with names and ``thou'') and LJ024 (two; the text of a political speech). This suggests that unusual vocabulary and oratorical delivery, rather than random clips, drive the largest errors; we did not isolate individual alignment mistakes.
\figbox{per_phone.png}{Test MAE per phoneme (best model).}{per_phone}

\subsection{HMM duration model versus real durations}\label{sec:eda-theory}
\figbox{hmm_vs_real.png}{Real duration histograms (bars) versus a geometric (single-state) and a 3-state negative-binomial model with the same mean.}{hmm}
Over the 12 most frequent phones the 3-state negative-binomial model fits \emph{better} only for longer phones (S $-2.87$ vs.\ $-3.35$ for geometric, Z $-2.66$ vs.\ $-3.13$, K $-2.62$ vs.\ $-3.15$; means 9--11 frames) and \emph{worse} for short ones (AH $-3.94$ vs.\ $-2.47$, DH $-4.64$ vs.\ $-2.15$; also N, D, R), because durations below 3 frames are impossible under it yet occur for 12.4\% of DH and 7.8\% of AH. Neither fixed shape fits every phone, which motivates a context-aware learned predictor. Caveat: MFA works in 10~ms frames while we count 11.6~ms frames, so the comparison is approximate.

% =====================================================================
\section{Discussion and Limitations}
\begin{itemize}
\item \textbf{Label noise.} Targets come from forced alignment, whose boundaries are themselves HMM estimates on a 10~ms grid; our predictors cannot beat the alignment's own error, and some bias of the HMM duration model is inherited by the labels.
\item \textbf{Irreducible variance.} Speaker choices (emphasis, rate within a sentence) are not recoverable from text, so $R^2$ has a ceiling below 1.
\item \textbf{Quantisation.} Converting 10~ms alignment frames to 11.6~ms mel frames adds rounding of up to half a frame.
\item \textbf{Single speaker, English, read speech.} Results may not transfer to expressive or multi-speaker data.
\item \textbf{Exclusions.} Removing \pctExcluded{} of utterances (mostly \texttt{spn} from abbreviations and proper names such as \emph{mr}, \emph{mrs}, \emph{fbi}) biases the retained set away from such text; extending the dictionary would recover most of them.
\item \textbf{No audio yet.} WER via a held-out ASR and the listening study require synthesis and belong to Phase~2.
\end{itemize}

\section{Conclusion and Phase~2 Outlook}
Phase~1 delivers validated duration labels (\nKept{} utterances), a frozen split and a documented comparison of duration predictors (best: \bestModel, MAE \maeBest{} frames).
Phase~2 will feed phonemes and (true or predicted) durations through a length regulator into a learned mel decoder (optionally diffusion-based, following Grad-TTS), render audio with a pretrained vocoder, and compare against duration-only synthesis using held-out-ASR WER, duration MAE and a small human intelligibility preference study, always on the same test split.

% =====================================================================
\section{Reproducibility}\label{sec:repro}
Public repository: \url{\repoURL}. Order of execution: notebooks 01 $\rightarrow$ 09 (see \texttt{README.md}). Shared settings live in \texttt{configs/data.yaml} (22{,}050~Hz, FFT 1024, hop 256, 80 mels); splits in \texttt{data/splits/}; every figure and table of this report is regenerated by notebook 09 from files in \texttt{results/}. Seeds: split 42, models 0. Software: Python, PyTorch, scikit-learn, MFA \mfaVersion. 

\section{Individual Contributions}\label{sec:contrib}
Each member owns one area and can explain every stage in the individual viva (Table~\ref{tab:contrib}).
\begin{table}[h]\centering\caption{Contributions.}\label{tab:contrib}
\footnotesize\begin{tabular}{@{}p{1.5cm}p{3.1cm}p{3.1cm}@{}}
\toprule Member & Main area & Notebooks / sections \\ \midrule
\teamA & Dataset exploration, EDA, theoretical rigor (HMM alignment, duration theory) & Notebooks 01--04; Sec.~\ref{sec:data}, \ref{sec:theory} \\
\teamB & Model application and rigor & Notebooks 06--08; Sec.~\ref{sec:models}, \ref{sec:results} \\
\teamC & Model application and rigor & Notebooks 06--08; Sec.~\ref{sec:models}, \ref{sec:results} \\
\teamD & Feature engineering, theoretical rigor (HMM alignment, duration theory)& Notebook 04--05; Sec.~\ref{sec:feat} \\
\teamE & Report, presentation, literature review, README & Sec.~\ref{sec:lit}, introduction, discussion, conclusion \\
\bottomrule
\end{tabular}
\end{table}

% =====================================================================
\begin{thebibliography}{99}
\bibitem{mcauliffe2017} M.~McAuliffe, M.~Socolof, S.~Mihuc, M.~Wagner, and M.~Sonderegger, ``Montreal Forced Aligner: Trainable text-speech alignment using Kaldi,'' in \emph{Proc. Interspeech}, 2017, pp.\ 498--502.
\bibitem{popov2021} V.~Popov, I.~Vovk, V.~Gogoryan, T.~Sadekova, and M.~Kudinov, ``Grad-TTS: A diffusion probabilistic model for text-to-speech,'' in \emph{Proc. ICML}, PMLR 139, 2021, pp.\ 8599--8608.
\bibitem{ljspeech2017} K.~Ito and L.~Johnson, ``The LJ Speech Dataset,'' 2017. [Online]. Available: \url{https://keithito.com/LJ-Speech-Dataset/}
\bibitem{rabiner1989} L.~R.~Rabiner, ``A tutorial on hidden Markov models and selected applications in speech recognition,'' \emph{Proc. IEEE}, vol.~77, no.~2, pp.\ 257--286, 1989.
\bibitem{viterbi1967} A.~J.~Viterbi, ``Error bounds for convolutional codes and an asymptotically optimum decoding algorithm,'' \emph{IEEE Trans. Inf. Theory}, vol.~13, no.~2, pp.\ 260--269, 1967.
\bibitem{ren2021} Y.~Ren \emph{et al.}, ``FastSpeech 2: Fast and high-quality end-to-end text to speech,'' in \emph{Proc. ICLR}, 2021.
\bibitem{kong2020} J.~Kong, J.~Kim, and J.~Bae, ``HiFi-GAN: Generative adversarial networks for efficient and high fidelity speech synthesis,'' in \emph{Proc. NeurIPS}, 2020.
\bibitem{friedman2001} J.~H.~Friedman, ``Greedy function approximation: A gradient boosting machine,'' \emph{Ann. Statist.}, vol.~29, no.~5, pp.\ 1189--1232, 2001.
\bibitem{hochreiter1997} S.~Hochreiter and J.~Schmidhuber, ``Long short-term memory,'' \emph{Neural Comput.}, vol.~9, no.~8, pp.\ 1735--1780, 1997.


\end{thebibliography}
\end{document}